\begin{figure}[t]
\centering
\begin{tikzpicture}
\begin{loglogaxis}[
    width=\columnwidth,
    height=6.5cm,
    xlabel={$t_{\max}$ (zile)},
    ylabel={Wall-time median al solverului (ms)},
    xmin=80, xmax=12000,
    ymin=0.5, ymax=2000,
    grid=both, grid style={line width=.1pt, draw=gray!20},
    minor tick num=9,
    legend style={at={(0.02, 0.98)}, anchor=north west, font=\footnotesize, fill=white, draw=gray!50},
    legend cell align={left},
    tick label style={font=\footnotesize},
    label style={font=\small},
    log basis x=10,
    log basis y=10,
]
\addplot[mark=square*, thick, color=blue!80!black, mark size=2.5pt]
    table[col sep=comma, x=t_max, y=julia] {data/scaling_wall_sidarthe.csv};
\addlegendentry{Julia (\texttt{Tsit5})}

\addplot[mark=triangle*, thick, color=green!50!black, mark size=2.5pt]
    table[col sep=comma, x=t_max, y=python] {data/scaling_wall_sidarthe.csv};
\addlegendentry{Python (\texttt{solve\_ivp} RK45)}

\addplot[mark=*, thick, color=red!80!black, mark size=2.5pt]
    table[col sep=comma, x=t_max, y=r] {data/scaling_wall_sidarthe.csv};
\addlegendentry{R (\texttt{deSolve} RK4)}

\end{loglogaxis}
\end{tikzpicture}
\caption{Scalarea wall-time-ului median al solverului SIDARTHE in functie de
orizontul $t_{\max}$, pentru cele trei limbaje. R prezinta o crestere
aproape liniara cu $t_{\max}$ ($\sim$$130\times$ intre 90 si 10\,000 de pasi),
dominata de inserarile rand-cu-rand in baza de date prin
\texttt{dbExecute()}; integrarea propriu-zisa cu \texttt{deSolve::ode} este
sub-secundara. Python ramane stabil ($4$--$8$~ms), iar Julia este mai rapid
in regimul amortizat dupa compilarea JIT, atingand $6$~ms la $10\,000$ de pasi.}
\label{fig:scaling_wall_time}
\end{figure}
